\section{Second-Order Systems and Bode Plot}

\subsection{Second-Order Standard Form}

\subsubsection{Second-order denominator}
\[
Q(s)=s^2+2\zeta\omega_n s+\omega_n^2.
\]
\(\omega_n\) is the natural angular frequency and \(\zeta\) is the damping factor.

\subsubsection{Second-order poles}
\[
p_{1,2}=-\zeta\omega_n\pm j\omega_n\sqrt{1-\zeta^2}.
\]
For \(0<\zeta<1\), the poles are complex conjugates with damped angular frequency \(\omega_d=\omega_n\sqrt{1-\zeta^2}\).

\subsubsection{Damping cases}
\[
\begin{array}{c|c}
\zeta & \text{case}\\ \hline
0<\zeta<1 & \text{underdamped}\\
\zeta=1 & \text{critically damped}\\
\zeta>1 & \text{overdamped}
\end{array}
\]
The damping factor determines whether the natural response oscillates, has repeated real poles, or has distinct real poles.

\subsubsection{Quality factor}
\[
Q=\frac{1}{2\zeta}.
\]
The quality factor is large for lightly damped second-order systems and small for strongly damped systems.

\subsubsection{Second-order lowpass, highpass, and bandpass}
\[
H_{\mathrm{LP}}(s)=K\frac{\omega_n^2}{s^2+2\zeta\omega_n s+\omega_n^2},
\]
\[
H_{\mathrm{HP}}(s)=K\frac{s^2}{s^2+2\zeta\omega_n s+\omega_n^2},
\qquad
H_{\mathrm{BP}}(s)=K\frac{2\zeta\omega_n s}{s^2+2\zeta\omega_n s+\omega_n^2}.
\]
The numerator determines whether low, high, or middle frequencies are passed.

\subsection{Bode Plot}

\subsubsection{Bode magnitude and phase}
\[
M_{\mathrm{dB}}(\omega)=20\log_{10}|H(j\omega)|,\qquad
\phi(\omega)=\arg H(j\omega).
\]
Bode plots show transfer-function magnitude in decibel and phase versus logarithmic frequency.

\subsubsection{Constant gain Bode rule}
\[
K \quad \Rightarrow \quad 20\log_{10}|K| \ \mathrm{dB},\qquad
\angle K=
\begin{cases}
0^\circ,&K>0,\\
180^\circ,&K<0.
\end{cases}
\]
Constant factors shift the magnitude plot without changing slope.

\subsubsection{Real zero and real pole Bode rules}
\[
1+\frac{s}{\omega_z}: \quad +20\ \mathrm{dB/dec}\ \text{after }\omega_z,
\]
\[
\frac{1}{1+s/\omega_p}: \quad -20\ \mathrm{dB/dec}\ \text{after }\omega_p.
\]
A first-order zero raises the asymptotic slope by \(20\) dB/decade. A first-order pole lowers it by \(20\) dB/decade.

\subsubsection{Origin pole and origin zero Bode rules}
\[
s^m \Rightarrow +20m\ \mathrm{dB/dec},\qquad
\frac{1}{s^m}\Rightarrow -20m\ \mathrm{dB/dec}.
\]
Zeros or poles at the origin set the initial Bode magnitude slope and add constant phase of \(\pm 90^\circ\) per order.

\subsubsection{Second-order pole pair Bode rule}
\[
\frac{\omega_n^2}{s^2+2\zeta\omega_n s+\omega_n^2}
\quad \Rightarrow \quad -40\ \mathrm{dB/dec}\ \text{after }\omega_n.
\]
A second-order pole pair contributes a total phase transition of approximately \(-180^\circ\). Resonance becomes pronounced for small \(\zeta\).

\subsubsection{Resonance peak}
\[
M_r=\frac{1}{2\zeta\sqrt{1-\zeta^2}},
\qquad
\omega_r=\omega_n\sqrt{1-2\zeta^2},
\quad 0<\zeta<\frac{1}{\sqrt{2}}.
\]
Underdamped lowpass systems have a resonance peak only for sufficiently small damping.
